\subsection{INDUCTIVE SETS}

DEFINITION 3. An ordered set E is said to be inductive if every totally ordered subset of E has an upper bound in E.

\minorheading{Examples}

\begin{enumerate}
\def\labelenumi{(\arabic{enumi})}
\item
  Let \(\mathfrak{F}\) be a set of subsets of a set \(A\) , ordered by inclusion, and such that for every totally ordered subset \(\mathfrak{G}\) of \(\mathfrak{F}\) the union of the sets of \(\mathfrak{G}\) belongs to \(\mathfrak{F}\) . Then \(\mathfrak{F}\) is inductive with respect to the relation \(\subset\) because the union of the sets of \(\mathfrak{G}\) is the least upper bound of \(\mathfrak{G}\) in \(\mathfrak{P}(A)\) .
\item
  An important example of a set of subsets which is inductive with respect to the relation \(\subset\) is the set \(\mathfrak{F}\) of graphs of mappings of subsets of a set A into a set B. For \(\mathfrak{F}\) is a subset of \(\mathfrak{P}(A \times B)\) , and to say that a subset \(\mathfrak{G}\) of \(\mathfrak{F}\) is totally ordered by inclusion means that the elements of \(\mathfrak{G}\) are graphs of mappings such that, given any two of these mappings, one is an extension of the other. It follows immediately that the union of the sets of \(\mathfrak{G}\) is an element of \(\mathfrak{F}\) (Chapter II, § 4, no. 6, Proposition 7). Hence the set \(\Phi(A, B)\) of mappings of subsets of A into B is inductive with respect to the order relation `` \(v\) extends \(u\) '' between \(u\) and \(v\) .
\end{enumerate}

* (3) It follows from the axiom of infinity (§ 6, no. 1) that the well-ordered set of natural integers is not inductive with respect to the relation ≤. *

THEOREM 2 (``Zorn's lemma''). Every inductive ordered set has a maximal element.

This theorem is a particular case of the following result :

PROPOSITION 4. Let E be an ordered set in which every well-ordered subset is bounded above; then E has a maximal element.

We shall say that an element \(v \in E\) is a strict upper bound of a subset \(X\) of \(E\) if \(v\) is an upper bound of \(X\) and \(v \notin X\) . Let \(\mathfrak{S}\) be the set of subsets of \(E\) which have a strict upper bound in \(E\) , and for each \(S \in \mathfrak{S}\) put \(p(S) = \tau_v(v\) is a strict upper bound of \(S)\) ; then \(p(S)\) is a strict upper bound of \(S\) . Applying Lemma 3 of no. 3 to \(\mathfrak{S}\) and \(p\) , we see that there exists a subset \(M\) of \(E\) and a well-ordering \(\Gamma\) on \(M\) which satisfies the conditions of Lemma 3; in particular, M has no strict upper bound in E. Furthermore, the ordering \(\Gamma\) is identical with that induced on M by the ordering on E. For in M the relation `` \(y \in \Gamma\langle x \rangle\) and \(x \neq y\) '' is equivalent to \(x \in S_y\) ; and since \(p(S_y) = y\) is an upper bound of \(S_y\) (with respect to the ordering on E), it implies \(x < y\) in E. But this means that the injection of M into E is a strictly increasing mapping (M being endowed with the ordering \(\Gamma\) ); and since M is totally ordered, it follows that the relations \(y \in \Gamma\langle x \rangle\) and \(x \leqslant y\) are equivalent in M ( \(\S 1\) , no. 12, Proposition 11). This being so, there exists by hypothesis an upper bound \(m\) of M in E; but since M has no strict upper bound, it follows that \(m\) is a maximal element of E.

COROLLARY 1. Let E be an inductive ordered set and let a be an element of E. Then there exists a maximal element m of E such that \(m \geqslant a\) .

For it follows from Definition 3 that the set F of elements \(x \geqslant a\) in E is inductive, and a maximal element of F is also maximal in E.

COROLLARY 2. Let \(\mathfrak{F}\) be a set of subsets of a set E such that, for every subset \(\mathfrak{G}\) of \(\mathfrak{F}\) which is totally ordered by inclusion, the union (resp. intersection) of the sets of \(\mathfrak{G}\) belongs to \(\mathfrak{F}\) ; then \(\mathfrak{F}\) has a maximal (resp. minimal) element.

